% Options for packages loaded elsewhere
\PassOptionsToPackage{unicode}{hyperref}
\PassOptionsToPackage{hyphens}{url}
%
\documentclass[
]{article}
\usepackage{amsmath,amssymb}
\usepackage{iftex}
\ifPDFTeX
  \usepackage[T1]{fontenc}
  \usepackage[utf8]{inputenc}
  \usepackage{textcomp} % provide euro and other symbols
\else % if luatex or xetex
  \usepackage{unicode-math} % this also loads fontspec
  \defaultfontfeatures{Scale=MatchLowercase}
  \defaultfontfeatures[\rmfamily]{Ligatures=TeX,Scale=1}
\fi
\usepackage{lmodern}
\ifPDFTeX\else
  % xetex/luatex font selection
\fi
% Use upquote if available, for straight quotes in verbatim environments
\IfFileExists{upquote.sty}{\usepackage{upquote}}{}
\IfFileExists{microtype.sty}{% use microtype if available
  \usepackage[]{microtype}
  \UseMicrotypeSet[protrusion]{basicmath} % disable protrusion for tt fonts
}{}
\makeatletter
\@ifundefined{KOMAClassName}{% if non-KOMA class
  \IfFileExists{parskip.sty}{%
    \usepackage{parskip}
  }{% else
    \setlength{\parindent}{0pt}
    \setlength{\parskip}{6pt plus 2pt minus 1pt}}
}{% if KOMA class
  \KOMAoptions{parskip=half}}
\makeatother
\usepackage{xcolor}
\usepackage[margin=1cm]{geometry}
\usepackage{graphicx}
\makeatletter
\def\maxwidth{\ifdim\Gin@nat@width>\linewidth\linewidth\else\Gin@nat@width\fi}
\def\maxheight{\ifdim\Gin@nat@height>\textheight\textheight\else\Gin@nat@height\fi}
\makeatother
% Scale images if necessary, so that they will not overflow the page
% margins by default, and it is still possible to overwrite the defaults
% using explicit options in \includegraphics[width, height, ...]{}
\setkeys{Gin}{width=\maxwidth,height=\maxheight,keepaspectratio}
% Set default figure placement to htbp
\makeatletter
\def\fps@figure{htbp}
\makeatother
\setlength{\emergencystretch}{3em} % prevent overfull lines
\providecommand{\tightlist}{%
  \setlength{\itemsep}{0pt}\setlength{\parskip}{0pt}}
\setcounter{secnumdepth}{-\maxdimen} % remove section numbering
\usepackage{amsmath}
\usepackage{mathtools}
\usepackage{amsthm}
\usepackage{color}
\usepackage{xcolor}
\usepackage{geometry}
\usepackage{enumerate}
\ifLuaTeX
  \usepackage{selnolig}  % disable illegal ligatures
\fi
\IfFileExists{bookmark.sty}{\usepackage{bookmark}}{\usepackage{hyperref}}
\IfFileExists{xurl.sty}{\usepackage{xurl}}{} % add URL line breaks if available
\urlstyle{same}
\hypersetup{
  pdftitle={Regression Analysis - STAT 482 - Exam 2: Due Dec 14, 2021},
  pdfauthor={Alex Towell (atowell@siue.edu)},
  hidelinks,
  pdfcreator={LaTeX via pandoc}}

\title{Regression Analysis - STAT 482 - Exam 2: Due Dec 14, 2021}
\author{Alex Towell
(\href{mailto:atowell@siue.edu}{\nolinkurl{atowell@siue.edu}})}
\date{}

\begin{document}
\maketitle

\newcommand{\sos}[1]{\mathrm{SS_{#1}}}
\newcommand{\ms}[1]{\mathrm{MS_{#1}}}
\newcommand{\sd}{\operatorname{sd}}
\newcommand{\var}{\operatorname{var}}
\newcommand{\expect}{\operatorname{E}}
\newcommand{\corr}{\operatorname{cor}}
\newcommand{\cov}{\operatorname{cov}}
\newcommand{\se}{\operatorname{se}}
\newcommand{\eval}[2]{\left. #1 \right\vert_{#2}}
\newcommand{\degf}[1]{\mathrm{df_{#1}}}
\newcommand{\entropy}{\operatorname{H}}
\newcommand{\param}[1]{\textcolor{blue}{\texttt{#1}}}
\newcommand{\ci}{\operatorname{CI}}

\hypertarget{problem-1}{%
\section{Problem 1}\label{problem-1}}

\begin{quote}
A beer distributor is interested in the amount of time to service its
retail outlets. Two factors are thought to influence the delivery time
(\(y\)) in minutes: the number of cases delivered (\(x_1\)) and the
distance traveled (\(x_2\)) in miles. A random sample of delivery time
data has been collected. The data is available on Blackboard as a csv
file.
\end{quote}

\vspace{50pt}

\hypertarget{a}{%
\subsection{(a)}\label{a}}

\begin{quote}
Provide an interpretation of a regression coefficient in a multiple
regression model.
\end{quote}

\vspace{50pt}

\hypertarget{b}{%
\subsection{(b)}\label{b}}

\begin{quote}
Compute \(b_1,b_2\), the estimated the regression coefficients for the
delivery time data.
\end{quote}

\vspace{50pt}

\hypertarget{c}{%
\subsection{(c)}\label{c}}

\begin{quote}
Compute \(t\) statistics for testing the effect of each input variable.
Explain what type of effect is being tested here.
\end{quote}

\vspace{50pt}

\hypertarget{d}{%
\subsection{(d)}\label{d}}

\begin{quote}
Compute \(\sos{R}(X_1)\) and \(\sos{R}(X_2|X_1)\). Explain what each sum
of squares represents.
\end{quote}

\vspace{50pt}

\hypertarget{e}{%
\subsection{(e)}\label{e}}

\begin{quote}
Test for a marginal effect of \(x_2\) against a model which includes no
other input variables. (Compute the test statistic and \(p\)-value.)
Provide an interpretation of the result, stated in the context of the
problem.
\end{quote}

\vspace{50pt}

\hypertarget{f}{%
\subsection{(f)}\label{f}}

\begin{quote}
Test for a partial effect of \(x_2\) against a model which includes
\(x_1\). (Compute the test statistic and \(p\)-value.) Provide an
interpretation of the result, stated in the context of the problem.
\end{quote}

\vspace{50pt}

\hypertarget{g}{%
\subsection{(g)}\label{g}}

\begin{quote}
Compute the correlation matrix. What feature of multidimensional
modeling is illustrated in this problem?
\end{quote}

\vspace{50pt}

\hypertarget{problem-2}{%
\section{Problem 2}\label{problem-2}}

\begin{quote}
A bakery is interested in the best formulation for a new product. A
small-scale experiment is conducted to investigate the relationship
between the product satisfaction (\(y\)), and the moisture content
(input 1) and sweetness (input 2) of the product. The input variables
have been coded \((x_1,x_2)\) for ease of calculation. The data is
available on Blackboard as a csv file.
\end{quote}

\vspace{50pt}

\hypertarget{a-1}{%
\subsection{(a)}\label{a-1}}

\begin{quote}
Provide a definition for an orthogonal design. Discuss an advantage to
using an orthogonal design.
\end{quote}

\vspace{50pt}

\hypertarget{b-1}{%
\subsection{(b)}\label{b-1}}

\begin{quote}
Provide a definition for an interaction effect.
\end{quote}

\vspace{50pt}

\hypertarget{c-1}{%
\subsection{(c)}\label{c-1}}

\begin{quote}
Fit an interaction model using the coded variables. Compute the
regression coefficient estimates and their standard errors.
\end{quote}

\vspace{50pt}

\hypertarget{d-1}{%
\subsection{(d)}\label{d-1}}

\begin{quote}
Write the estimated regression as a function of \(x_1\) for
\(x_2 = 1,0,-1\).
\end{quote}

\vspace{50pt}

\hypertarget{e-1}{%
\subsection{(e)}\label{e-1}}

\begin{quote}
Create interaction plots for both the interaction model and the additive
effects model.
\end{quote}

\vspace{50pt}

\hypertarget{f-1}{%
\subsection{(f)}\label{f-1}}

\begin{quote}
Test for an interaction effect. (Compute the test statistic and
\(p\)-value.)
\end{quote}

\vspace{50pt}

\hypertarget{problem-3}{%
\section{Problem 3}\label{problem-3}}

\begin{quote}
An engineer is interested in comparing three chemical processes
(categorical input with groups \(A\),\(B\),\(C\)) for manufacturing a
compound. It is suspected that the impurity (continuous input \(x\)) of
the raw material will affect the yield (response variable \(y\)) of the
product. The data is available on Blackboard as a csv file.
\end{quote}

\vspace{50pt}

\hypertarget{a-2}{%
\subsection{(a)}\label{a-2}}

\begin{quote}
Define indicator variables \(I_1\) and \(I_2\) using chemical process
\(C\) as the baseline level.
\end{quote}

\vspace{50pt}

\hypertarget{b-2}{%
\subsection{(b)}\label{b-2}}

\begin{quote}
Write an additive model for response \(y\) using continuous input
variable \(x\) and indicator variables \(I_1\), \(I_2\).
\end{quote}

\vspace{50pt}

\hypertarget{c-2}{%
\subsection{(c)}\label{c-2}}

\begin{quote}
Write a regression function for each of the chemical processes.
\end{quote}

\vspace{50pt}

\hypertarget{d-2}{%
\subsection{(d)}\label{d-2}}

\begin{quote}
Provide an interpretation for each effect parameter, stated in the
context of the problem.
\end{quote}

\vspace{50pt}

\hypertarget{e-2}{%
\subsection{(e)}\label{e-2}}

\begin{quote}
Compute interval estimates for each of the effect parameters.
\end{quote}

\vspace{50pt}

\hypertarget{f-2}{%
\subsection{(f)}\label{f-2}}

\begin{quote}
Create a scatterplot of the data with the estimated regression lines.
\end{quote}

\end{document}
